\documentclass[11pt]{article}

\usepackage{amsmath, amssymb, amsthm}
\usepackage{graphicx}
\usepackage{bm}
\usepackage{physics}
\usepackage{geometry}
\usepackage{hyperref}
\usepackage{caption}
\usepackage{color}

\geometry{margin=1in}

\title{Constraint Torques, Wrist Joint Mechanics, and Force Transmission in the Golf Swing:\\
A Unified Treatment from Robotics and Biomechanics}
\author{ }
\date{ }

\begin{document}
\maketitle

\begin{abstract}
The golf swing is a dynamic, highly coupled motion in which mechanical constraints, passive torques, 
and nonlinear interactions play a decisive---yet often underappreciated---role.  
This article develops a rigorous treatment of \emph{constraint forces and torques} in multibody systems 
and applies the framework to the left wrist, modeled as a universal joint connecting the hand to the 
forearm.  

We explore (i) how such torques arise, (ii) how they are computed in modern robotics and in 
Simscape/Simulink, (iii) how they produce effective torque transmission even though they do no mechanical 
work, and (iv) how variations in grip orientation may modulate both swing speed and clubface variability.  

The core hypothesis examined here is that the \emph{orientation of the club within the hand relative to the wrist's universal-joint axes} 
changes the way constraint torques are routed through the wrist and into the shaft.  
Because different axes of the club have drastically different inertias, this alignment may influence both 
the magnitude of transmitted torques and their variability---helping to explain practical coaching advice, 
such as ``grip in the fingers, not the palm.''  
\end{abstract}

\section{Introduction}

Golf instruction has long emphasized seemingly simple cues such as 
``hold the club in the fingers,'' 
``avoid a palm grip,'' and 
``keep the club and arms slightly offset in plane.''  
These rules-of-thumb typically lack a rigorous mechanical foundation despite their empirical support.  
This paper argues that a substantial part of these phenomena emerges naturally once the wrist 
is treated as a \emph{universal joint} transmitting both active torques and 
\emph{constraint torques}---passive, nonlinear, configuration-dependent torques that arise 
when motion along forbidden directions must be prevented.

Although widely recognized in robotics and mechanical engineering, constraint torques have been 
largely overlooked in the biomechanics of striking motions.  
This is unfortunate because the golf swing is precisely the sort of fast, multi-DOF, highly coupled 
task where constraints dominate the internal force landscape.  

The left wrist---a mechanical bottleneck between the powerful proximal muscles of the body and the 
clubhead---is particularly sensitive to these interactions.  
During the downswing and especially near impact, the wrist transmits large forces and torques, many of 
which are \emph{not} generated directly by the golfer’s muscles.  
Instead, they arise passively from the coupling between:
\begin{itemize}
    \item the universal joint geometry of the wrist,
    \item the rotation of the forearm,
    \item the inertia of the club,
    \item and the nonholonomic constraints that forbid rotation about the wrist’s third axis.
\end{itemize}

As we will show, these constraint torques can be large, highly variable, and extremely influential 
in determining club dynamics.

\section{The Wrist as a Universal Joint}

The human wrist is anatomically complex, but for the purpose of dynamic modeling, a universal-joint 
approximation is remarkably effective.  
A universal joint provides two rotational degrees of freedom arranged orthogonally:
\[
(\theta_x,\; \theta_y)
\]
corresponding approximately to flexion/extension and radial/ulnar deviation.  
A third rotation (forearm pronation--supination) occurs in the radius–ulna pair, not the wrist itself.

Thus the left arm–hand–club chain can be idealized as:
\[
\text{Forearm} \quad 
\overset{\text{(pronation DOF)}}{\longrightarrow} \quad
\text{Radial/Ulnar Axis} \quad
\overset{\text{(ulnar/radial)}}{\longrightarrow} \quad
\text{Flex/Ext Axis} \quad
\longrightarrow \text{Hand/Club}
\]

Even this simple model produces surprisingly rich behavior.  
When the model is driven forward dynamically, torques sensed at the joint include:

\begin{itemize}
    \item the \textbf{applied torques} $\tau_x,\tau_y$, which are set by the golfer or simulation;
    \item the \textbf{constraint torque} $\tau_z$, a passive torque enforcing forbidden rotation about the third axis;
    \item \textbf{additional passive components} altering the sensed $\tau_x,\tau_y$, arising from inertial coupling, nonlinear constraints, and momentum transfer.
\end{itemize}

In other words, even in a nominally 2-DOF universal joint,
\[
\bm{\tau}_{\text{sensed}} = 
\begin{bmatrix}
\tau_x \\[3pt]
\tau_y \\[3pt]
\tau_z
\end{bmatrix}
\neq 
\begin{bmatrix}
\tau_{\text{input},x} \\[3pt]
\tau_{\text{input},y} \\[3pt]
0 
\end{bmatrix}.
\]

This mismatch---the source of much confusion in biomechanical interpretation---is central to our analysis.

\section{Where Constraint Torques Come From: A Robotics Perspective}

In a constrained multibody system, the dynamics take the standard form
\begin{equation}
    M(q)\ddot{q} + C(q,\dot{q}) + G(q)
    = \tau_{\text{act}} + J_c(q)^\top \lambda,
    \label{eq:dyn}
\end{equation}
where:
\begin{itemize}
    \item $M(q)$ is the generalized inertia matrix,
    \item $C(q,\dot{q})$ captures Coriolis/centrifugal terms,
    \item $G(q)$ is gravity,
    \item $\tau_{\text{act}}$ are \emph{applied} (muscle/driver) torques,
    \item $J_c$ is the constraint Jacobian expressing forbidden velocities,
    \item $\lambda$ are the \emph{constraint forces/torques}.
\end{itemize}

The constraint forces satisfy
\[
J_c(q)\dot{q} = 0 \quad \Rightarrow \quad J_c(q)\ddot{q} + \dot{J}_c(q,\dot{q})\dot{q} = 0.
\]

Solving this gives 
\[
\lambda = -(J_c M^{-1} J_c^\top)^{-1} \left(J_c M^{-1}(\tau_{\text{act}} - C - G) + \dot{J}_c \dot{q}\right).
\]

The constraint torques observed at a joint are:
\[
\tau_{\text{constraint}} = J_c^\top \lambda.
\]

These torques satisfy several crucial properties:
\begin{enumerate}
    \item \textbf{They do no work.}  
    Constraint torques act along directions of blocked motion.  
    Since the generalized velocity in that direction is zero,
    \[
    \tau_{\text{constraint}}^\top \dot{q} = 0.
    \]
    Nonetheless, they can be enormous in magnitude.
    
    \item \textbf{They passively transmit force and torque}.  
    They are the reason a universal joint can transmit torsion along the shaft axis.

    \item \textbf{They depend on system geometry, momentum, and applied torques}.  
    They are nonlinear, configuration-dependent, and prone to large fluctuations.

    \item \textbf{They can generate apparent torques about axes with no DOF}.  
    In a universal joint: rotating about both active axes simultaneously produces passive torque about the third.
\end{enumerate}

\section{Demonstration: Why Flexion + Ulnar Torque Produces a Passive Z-Torque}

Try this yourself:  
Hold your left wrist with your right hand to prevent forearm rotation.  
Now torque your left wrist simultaneously toward flexion and ulnar deviation.

You will feel a twisting torque about the forearm axis—even though you are not actively applying any.

This is the textbook universal-joint effect: combining rotations about both active axes 
forces the system into a configuration where a twisting torque is needed to maintain the 
constraint that no forearm rotation is allowed.

Simulink reproduces this exactly:
\[
\bm{\tau}_{\text{sensed}} = 
\begin{bmatrix}
\tau_x + \tau_{x,\text{passive}} \\
\tau_y + \tau_{y,\text{passive}} \\
\tau_{z,\text{constraint}}
\end{bmatrix}.
\]

\section{Why This Matters for Force Transmission to the Club}

Constraint torques are a gift from the mechanical universe:  
they allow the torso, shoulders, and forearm to deliver torque to the club around axes the wrist cannot rotate about.  

In the context of a golf swing:
\begin{itemize}
    \item Large body-generated torques and angular momenta arrive at the wrist.
    \item The wrist denies rotation about the third axis.
    \item The resulting constraint torque is routed into the club.
    \item Depending on the grip, this torque goes into either:
    \begin{itemize}
        \item \textbf{the in-plane axis} (high inertia, low effect on face angle), or
        \item \textbf{the shaft axis} (low inertia, strong effect on face angle).
    \end{itemize}
\end{itemize}

This routing is where grip geometry enters the picture.

\section{Grip Orientation as a Mechanical Routing Problem}

Let the club coordinate frame at the hands be defined by:
\[
(\hat{x}_{\text{club}}, \hat{y}_{\text{club}}, \hat{z}_{\text{club}}).
\]

Near impact:
\begin{itemize}
    \item $\hat{x}$ is roughly along the loft line,
    \item $\hat{y}$ approximately normal to the swing plane,
    \item $\hat{z}$ along the shaft.
\end{itemize}

The wrist's universal-joint “hand-side” axes will be some orientation:
\[
(\hat{X}_{\text{wrist}}, \hat{Y}_{\text{wrist}}, \hat{Z}_{\text{wrist}}).
\]

The constraint torque generated at the wrist is primarily about $\hat{Z}_{\text{wrist}}$.
Thus the effective torque delivered to the club is:
\[
\bm{\tau}_{\text{club}} = R_{\text{wrist}\to\text{club}} \, \bm{\tau}_{\text{constraint}}.
\]

If the grip is in the \textbf{fingers}, $\hat{Z}_{\text{wrist}}$ is closer to the plane normal $\hat{y}_{\text{club}}$:
\begin{itemize}
    \item torque is routed into a high-inertia axis,
    \item clubface orientation becomes \emph{less sensitive},
    \item translational acceleration of the clubhead is favored (good for speed).
\end{itemize}

If the grip is \textbf{deep in the palm}, the wrist's $Z$-axis aligns closer with the shaft:
\begin{itemize}
    \item torque is routed into the low-inertia shaft-axis rotation,
    \item face angle becomes \emph{more sensitive} and more variable,
    \item less torque contributes to translating the clubhead (bad for speed).
\end{itemize}

This is a purely mechanical argument supporting the long-standing instruction tradition.

\section{Why Constraint Torques Are Variable (and Why That Matters)}

Constraint torques depend on:
\begin{itemize}
    \item instantaneous joint configuration,
    \item forearm rotation rate,
    \item angular momentum of the club,
    \item coupling terms in $C(q,\dot{q})$,
    \item and external torques from proximal segments.
\end{itemize}

Because all of these change rapidly during the downswing, constraint torques tend to be:
\[
\text{large, oscillatory, nonlinear, and unpredictable.}
\]

This variability inevitably affects:
\begin{itemize}
    \item clubface rotation,
    \item wrist load tolerance,
    \item torque transmission efficiency,
    \item consistency of delivery.
\end{itemize}

Thus grip geometry may be an important modulator of complexity in the wrist’s control problem.

\section{The Kinematics–to–Kinetics Challenge in Simulation}

In Simscape, when the joint is driven kinematically, the software internally solves 
Equation~\eqref{eq:dyn} for $\lambda$ and reports:
\[
\tau_{\text{sensed}} = \tau_{\text{act}} + \tau_{\text{constraint}}.
\]

Attempting to invert this mapping to obtain the ``actuator torques needed to reproduce a motion’’ 
is nontrivial because:
\begin{itemize}
    \item constraint torques depend on \emph{both} $\tau_{\text{act}}$ and $\ddot{q}$,
    \item the wrist is dynamically coupled to the entire upper body,
    \item forearm rotation couples strongly into constraint torques,
    \item the mapping is nonlinear and configuration-dependent.
\end{itemize}

This is the same difficulty addressed in modern robotics textbooks, and one reason 
you found Northwestern’s Modern Robotics resources so valuable—they explicitly cover constrained dynamics 
and the role of constraint forces, whereas most biomechanics texts do not.

\section{Hypothesis: Grip Orientation Influences Face-Angle Variability}

We restate your hypothesis in mechanical form:

\begin{quote}
\textbf{Hypothesis.}  
The orientation of the club in the left hand determines the mapping from wrist constraint torque 
to club angular acceleration.  
Finger-based grips align the wrist’s constraint-torque axis with a high-inertia axis of the club, 
reducing face-angle sensitivity and enhancing translation of the clubhead, while palm-based grips align 
constraint torques with the shaft axis, increasing face-angle variability and reducing clubhead speed.
\end{quote}

This is testable both numerically and experimentally.

\section{Proposed Simulation Experiments}

\begin{enumerate}
    \item Drive a universal-joint wrist with measured $\theta_x(t),\theta_y(t)$ kinematics.
    \item Vary the grip rotation $R_{\text{wrist}\to\text{club}}$ across realistic ranges 
    (e.g., finger vs palm grips).
    \item Extract:
    \begin{itemize}
        \item sensed constraint torques,
        \item resulting clubface angular accelerations,
        \item translational clubhead speed,
        \item statistical variability under small perturbations (noise injection).
    \end{itemize}
    \item Compare high- vs low-inertia routing conditions.
\end{enumerate}

This will quantitatively evaluate whether the traditional golf-instruction insights have mechanical roots.

\section{Conclusion}

Constraint torques play a decisive role in the dynamics of the golf swing.  
The universal-joint nature of the wrist, combined with strong inertial and constraint couplings, 
means that:
\begin{itemize}
    \item large torques reach the club without being actively produced at the wrist,
    \item these torques can be routed differently depending on grip geometry,
    \item and this routing materially affects both clubhead speed and clubface variability.
\end{itemize}

The mechanical case for preferring a ``finger grip’’ over a ``palm grip’’ emerges naturally 
from this analysis: finger-based grips align passive torque transmission with a high-inertia axis, 
reducing unwanted face rotation while increasing the contribution of constraint torques to clubhead translation.

This work opens a path for rigorous, simulation-backed mechanical understanding of golf swing variability, 
with potential applications in instruction, equipment design, and personalized fitting.
\end{document}
